% Generated by scripts/assemble.py; corrections belong in source/corrections.json.
\examhead{126}{中国海洋大学 2026 年数学分析试卷(回忆版)}{120}
% p121-b02; source PDF page 121
\blockanchor{p121-b02}
(注. 题目来自考生回忆, 部分表述或数据可能有误)\par
% p121-b03; source PDF page 121
\blockanchor{p121-b03}
\question{1}计算题.\par
% p121-b04; source PDF page 121
\blockanchor{p121-b04}
(1) 求极限 $\lim_{n\to \infty}\left(\ln n - \sum_{k = 1}^{n}\frac{\ln k}{n}\right)$.\par
% p121-b05; source PDF page 121
\blockanchor{p121-b05}
(2) 求极限 $\lim_{x\to 0}\frac{(1 + x)^x - \cos\frac{x}{2}}{x(\sin x - \sin\frac{x}{2})}.$\par
% p121-b06; source PDF page 121
\blockanchor{p121-b06}
(3) 求积分  $\int_{0}^{1}\frac{x^{3}-x}{\ln x}dx$ ，需说明计算的合理性.\par
% p121-b07; source PDF page 121
\blockanchor{p121-b07}
\question{2}判断题. 正确的给出证明, 错误的给出反例.\par
% p121-b08; source PDF page 121
\blockanchor{p121-b08}
(1) 任何数列都有单调的子列.\par
% p121-b09; source PDF page 121
\blockanchor{p121-b09}
(2) 若级数 $\sum_{n=1}^{\infty}\left(a_{3n-2} + a_{3n-1} + a_{3n}\right)$ 收敛, 且 $\lim_{n\to\infty}a_n = 0$, 则 $\sum_{n=1}^{\infty}a_n$ 收敛.\par
% p121-b10; source PDF page 121
\blockanchor{p121-b10}
(3) 设 $f(x)$ 在 $(- \infty, + \infty)$ 上二阶可导, 则对任意的 $\xi \in (-\infty, + \infty)$, 一定存在 $x_1, x_2 \in (-\infty, + \infty)$, 满足 $f'(\xi) = \frac{f(x_2) - f(x_1)}{x_2 - x_1}$.\par
% p121-b11; source PDF page 121
\blockanchor{p121-b11}
(4) 若 $f(x)$ 在 $[a, +\infty)$ 上单调非负且无界, 则 $\int_{a}^{+\infty} f(x) \, \mathrm{d}x$ 必发散.\par
% p121-b12; source PDF page 121
\blockanchor{p121-b12}
\question{3}解答如下问题:\par
% p121-b13; source PDF page 121
\blockanchor{p121-b13}
(1) (可能有误) 若 $\lim_{n\to \infty}\frac{a_1 + a_2 + \cdots + a_n}{n} = a,$ 则 $\lim_{n\to \infty}a_n = a.$\par
% p121-b14; source PDF page 121
\blockanchor{p121-b14}
(2) 若 $f(x)$ 在 $[a, +\infty)$ 上可导，且 $\lim_{x \to a^{+}} f(x) = \lim_{x \to +\infty} f(x) = A$ (A 为有限数)，则存在 $\xi \in (a, +\infty)$，使得 $f'(\xi) = 0$.\par
% p121-b15; source PDF page 121
\blockanchor{p121-b15}
\question{4}证明函数项级数 $\sum_{n=1}^{\infty} a_n \frac{1}{n!} \int_0^x e^{-t} t^n \, \mathrm{d}t$ 一致收敛的充要条件为级数 $\sum_{n=1}^{\infty} a_n$ 收敛.\par
% p121-b16; source PDF page 121
\blockanchor{p121-b16}
\question{5}设 $z = z(x, y)$ 二阶连续可微, 求 $a$ 的值, 使得通过 $u = x + 3y$, $v = x + ay$ 将 $6z_{xx} + z_{xy} - z_{yy} = 0$ 化为 $z_{uv} = 0$.\par
% p121-b17; source PDF page 121
\blockanchor{p121-b17}
\question{6}解答如下问题:\par
% p121-b18; source PDF page 121
\blockanchor{p121-b18}
(1) (可能有误)证明级数 $\sum_{n=1}^{\infty} \frac{(-1)^n}{x^n + \frac{2^n}{n(n+1)}}$ 在 $(- \infty, +\infty)$ 上一致收敛.\par
% p121-b19; source PDF page 121
\blockanchor{p121-b19}
(2) 求 $\lim_{x\to 0}\sum_{n = 1}^{\infty}\frac{(-1)^n}{x^n + \frac{2^n}{n(n + 1)}}.$\par
% p121-b20; source PDF page 121
\blockanchor{p121-b20}
\question{7}计算 $\iint_{D} \sqrt{|y - x^2|} \, \mathrm{d}x \, \mathrm{d}y$, 其中 $D = \{(x, y) \mid x^2 \leq y \leq 2\}$.\par
